\section{Diamond Interchange Zone Design}
\label{sec:diamond-design}

\subsection{Zone Definition and Scope}

A diamond interchange zone is a 50-mile corridor centered on a T1/T1 junction within which freight can enter or exit either T1 corridor at multiple independent points. The zone boundary is defined by access-controlled transfer eligibility: any point within the zone where a freight vehicle can legally and physically transition between the two T1 corridors at interstate grade standard is a zone connection point. The zone achieves k $\geq$ 3 when at least three such connection points exist, each on geometrically independent rights-of-way such that no single structural failure or closure event affects more than one.

The 50-mile radius is a policy parameter derived from freight economics. The maximum acceptable detour for a freight vehicle diverting from a primary T1 path to a zone alternate is approximately 50--75 additional route miles, corresponding to 45--70 minutes of additional transit time at 65 mph. At ATRI operating cost of \$225 per truck-hour \citep{ATRI_costs2024}, a 60-minute detour costs \$225 --- below the threshold at which a carrier would declare the diversion economically infeasible for a single-day disruption event. A 50-mile zone radius captures all connection points within this threshold.

\subsection{Key Zone Components}

\textbf{Component 1: Primary interchange node.} The primary interchange node is an elevated direct-connect ramp from T1 corridor A to T1 corridor B, grade-separated from all local traffic movements. Access is restricted to freight-classified vehicles (commercial motor vehicles with USDOT registration operating under FMCSA authority). Design speed is 65 mph. Curve radius minimum is 600 feet, adequate for 65-mph combination vehicle operation \citep{HCM7}. Maximum grade is 5\%; grades above 3\% require a dedicated truck climbing lane on the ascending approach. Pavement design is for 10-million-ESAL loading (25-year heavy commercial design life). Typical construction cost: \$150--200 million per primary node, depending on urban vs. rural context and extent of grade separation required.

\textbf{Component 2: Secondary ramp access points.} The zone requires 2--4 additional access points within the 50-mile zone, each providing a distinct edge-disjoint path between the two T1 corridors. Secondary ramps need not be fully grade-separated from local traffic, but must carry 80,000-lb legal loads and maintain 55-mph or greater design speed. Secondary ramps typically use existing US highway or T2 corridor rights-of-way upgraded to freight-adequate geometry. Each secondary ramp adds one edge-disjoint path to the zone graph; three secondary ramps achieve k = 4 when combined with the primary node. Per-ramp cost: \$30--50 million for geometry improvement and weight-limit upgrade; \$80--120 million for partial grade separation at the T1 junction points.

\textbf{Component 3: Weigh station bypass.} Transponder-equipped freight vehicles (enrolled in the zone's managed access program) bypass physical weigh stations within the zone on both T1 corridors. The bypass removes a chronically congested bottleneck at the primary interchange approach: on T1 corridors with 8,000--12,000 commercial vehicles per day, weigh station backup during peak periods adds 15--25 minutes to interchange approach time, degrading the effective speed of the primary connection point. Weigh station bypass infrastructure (in-motion weight sensors, transponder readers, bypass lanes) costs \$3--5 million per station location.

\textbf{Component 4: Geometric design standards.} All zone components share common geometric standards drawn from FHWA managed lane design guidance \citep{FHWA_ML2021} and HCM7 freeway standards \citep{HCM7}: minimum 600-foot horizontal curve radius at 65-mph design speed; maximum 5\% grade (with truck climbing lane above 3\%); minimum 12-foot travel lanes with 10-foot right shoulder and 4-foot left shoulder; pavement rated for 10-million-ESAL design life; overhead clearance minimum 16.5 feet for combination vehicle access.

\subsection{Design Speed Versus Cost Tradeoff}

The 65-mph design speed standard for primary zone nodes adds approximately 30\% to construction cost compared to a 45-mph interchange design. The premium is driven primarily by curve geometry: a 45-mph design uses a 300-foot minimum radius, while a 65-mph design requires 600 feet. On constrained urban sites, the 300-foot additional radius translates to additional right-of-way acquisition, retaining structure construction, or elevated structure length. The 65-mph standard is retained despite this premium because the zone's freight throughput benefit depends on sustained-speed operation: a direct-connect ramp that requires deceleration to 45 mph does not enable truck platooning and does not eliminate the speed differential that creates merge conflicts at the primary interchange \citep{Browand2004}.

\subsection{Case Studies: Three Priority Intersections}

\textbf{Atlanta: I-75/I-85.} The Atlanta metropolitan area presents an unusual geometry: I-75 and I-85 share a common 7.2-mile segment through downtown Atlanta (the Downtown Connector, I-75/I-85 multiplex), making the two T1 corridors effectively a single entity through the urban core. The diamond zone must therefore be defined around the I-285 perimeter ring, which provides the zone's primary alternate connectivity. Secondary ramps are located at I-285/I-20 (western quadrant) and I-285/I-16 (eastern approach toward Savannah). The combination of four independent path options --- the primary Downtown Connector, I-285 north segment, I-285 south segment, and the I-20/I-16 bypass --- achieves k = 4 for the zone. Estimated capital cost: \$380 million, including \$160 million for elevated direct-connect freight ramps at the I-285 north and south interchanges and \$220 million for geometry upgrades on I-285 freight approaches.

\textbf{Jacksonville: I-10/I-95.} Jacksonville presents a T-intersection geometry: I-10 terminates at I-95 at the eastern end of its transcontinental run. The single existing interchange carries all I-10 freight transitioning to I-95 northbound or southbound through a standard four-ramp diamond. A direct-connect freight ramp parallel to the existing interchange --- elevated, grade-separated from the local interchange movements, with dedicated I-95 N and I-95 S connections --- allows I-10 freight to join I-95 without entering the local interchange weave zone. Secondary ramps using I-295 (the Jacksonville beltway) and US-1 (upgraded to freight standard at the Buckman Bridge) provide two additional edge-disjoint paths. Estimated capital cost: \$210 million, reflecting the lower complexity of the T-intersection geometry compared to Atlanta's multiplex condition.

\textbf{Toledo: I-75/I-90.} Toledo presents the most geometrically challenging design in the three-priority group. I-75 runs north-south and I-90 runs east-west; the existing interchange is a partial cloverleaf with local access ramps integrated into the interchange footprint, creating the k = 1 SPF condition. Full diamond conversion requires complete grade separation of I-75 through-traffic from I-90 through-traffic, relocation of local access ramps to adjacent interchange locations, and construction of direct-connect freight ramps on all four turning movements (NB-EB, NB-WB, SB-EB, SB-WB). Secondary ramps using US-20 (upgraded to 80,000-lb standard) and I-280 (the Toledo beltway) provide additional zone paths. Local traffic relocation is the primary cost driver: realigning three existing local interchanges adds \$80--100 million to what would otherwise be a \$240 million primary-node project. Total estimated capital cost: \$340 million.

\subsection{Zone Performance Metric}

The operational target for a completed diamond interchange zone is a zone transfer time of 12 minutes or less: the elapsed time from a freight vehicle's decision point (last feasible diverge from the primary T1) to merge onto the alternate T1 corridor, measured at design speed on a zone with k $\geq$ 3 connectivity. The 12-minute target is derived from the freight economics constraint: a 12-minute transfer adds 0.2 hours of transit time at \$225/hr ATRI cost, or \$45 per vehicle --- an acceptable disruption premium for an unplanned diversion event. Zones exceeding 12-minute transfer time require either additional secondary ramp development within the zone or speed standard upgrades on existing zone connections.
